
%% Put this in main.tex and it will work

\documentclass{article}

% Load your entire custom package library with one line.
\usepackage{bsk_latex_lib/bsk_latex_lib}

% --- Document Configuration ---
% All high-level settings for this specific document are placed here.
\setupProtokoll{
  sektion = {boss},
  plats   = {Karlskrona/Microsoft Teams},
  datum   = {2025-mm-dd},
}

% --- Person Database ---
% Define all relevant people and their roles.
\newperson{\Ordfarande} { name = {Ordförande},    role = {Ordförande} }
\newperson{\Person}  { name = {Namn Namnson},     role = {Vice-ordförande} }
\newperson{\Namn} { name = {Namn Namnson},       role = {Sekreterare} }
% \newperson{\Namn} { name = {Namn Namnson}, role = {Ekonomiansvarig} }
% \newperson{\Namn} { name = {Namn Namnson},   role = {Ledamot} }
% \newperson{\Namn}  { name = {Namn Namnson},   role = {Ledamot} }
% \newperson{\Namn} { name = {Namn Namnson}, role = {Ledamot} }
% \newperson{\Namn} { name = {Namn Namnson},    role = {Suppleant} }
% \newperson{\Namn} { name = {Namn Namnson},   role = {Suppleant} }

\begin{document}

% --- Meeting-Specific Data ---
% This is the only section that needs to be filled in for each new meeting.

% 1. Clear any data from previous meetings (good practice).
\clearAttendance

% 2. Assign the officials for THIS meeting.
\setMeetingOfficials{\Ordfarande}{\Person}{\Namn} % Chairperson, Secretary, Adjuster

% 3. Set the attendance for THIS meeting.
\addPresent{\Ordfarande}
\addPresent{\Person}
\addPresent{\Namn}

% \addAbsent{\Namn}
% \addAbsent{\Namn}

% --- Automated Document Generation ---

% This single command now generates the attendance lists and all
\printFormalia

% --- Meeting Content (Ärenden) ---
% The unique content of the meeting is written here.

\paragraf{Skrivelser}{
  Mail från xx angående yy \\
  Mail från xx angående yy \\
}

\paragraf{Rapporter}{
  % UPDATED: Use \person{\Melker} to get the name
  \person{\Ordfarande} har:
  \begin{itemize}
    \item Inget att rapportera.
  \end{itemize}
  \person{\Namn} har:
  \begin{itemize}
    \item Inget att rapportera.
  \end{itemize}
}

\paragraf{Tidigare action-points}{}

\paragraf{Studiebevakning}{
  \textbf{Möte som vi kanske varit på}\\
  Brödtext
  \enter

  \textbf{Annat möte eller punkt}\\
  Brödtext
}
\paragraf{Feedback-soffan}{
  Inget att rapportera.
}

\paragraf{Studiesocialt}{
  \textbf{Event}\\
  Brödtext
  \enter

  \textbf{Annat event}\\
  Brödtext
}

\paragraf{Styrelsen}{}
\paragraf{Vibe check}{}

% --- Closing Section ---
% The \printClosing command handles the decision summary and signatures.
% The paragraphs for Övrigt, Kalender, etc., are placed before it.

\paragraf{Övrigt}{}
\paragraf{Kalender}{
  \begin{tabular}{ll}
  Nästa styrelse möte & yyyy-mm-dd\\
  Event 1 & yyyy-mm-dd\\
  Event 2 & yyyy-mm-dd\\
  Möte 1 & yyyy-mm-dd\\
  Möte 2 & yyyy-mm-dd\\
  \end{tabular}
}

% §15 for Beslutssammanfattning is now generated automatically by \printClosing.

\paragraf{Action points}{
  \person{\Ordfarande} ska:
  \begin{itemize}
    \item göra X
    \item göra Y
  \end{itemize}
  \person{\Namn} ska:
  \begin{itemize}
    \item Maila X
    \item göra klart Y
  \end{itemize}
}


\printClosing

\end{document}
